\documentclass[letterpaper]{article}

\usepackage{natbib,alifeconf}  %% The order is important
\usepackage{graphicx}
\usepackage{amsthm}
\usepackage{amsmath}
\usepackage{booktabs}
\usepackage{amssymb}
\usepackage{wasysym}
\usepackage{glossaries}
\usepackage{biocon}
\usepackage{todonotes}
\usepackage{standalone}
\usepackage{hyperref}

% my symbols %%%%%%%%%%%%%
\newcommand{\pitch}{\pi}
\newcommand{\yaw}{\epsilon}
\newcommand{\azimuth}{\alpha}
\newcommand{\elevation}{\eta}

% Read in data ##########
\input{variables}

% my acronyms %%%%%%%%%%%%%
\newacronym{uav}{UAV}{Unmanned-Aerial Vehicle}
\newcommand{\uav}{\gls{uav}}
\newcommand{\uavs}{\glspl{uav}}

\newacronym{catd}{CATD}{Constant Absolute Target Direction}
\newcommand{\catd}{\gls{catd}}

\newacronym{ait}{AIT}{Acoustic Image Theory}
\newcommand{\ait}{\gls{ait}}

\newacronym{pcb}{PCB}{printed circuit board}
\newcommand{\pcb}{\gls{pcb}}
\newcommand{\pcbs}{\glspl{pcb}}

\newacronym{hrtf}{HRTF}{head-related transfer function}
\newcommand{\hrtf}{\gls{hrtf}}
\newcommand{\hrtfs}{\glspl{hrtf}}

% my species %%%%%%%%%%%%%%

\newanimal{ca}{genus=Cordulegaster,epithet=annulatus}
\newanimal{mm}{genus=Micronycteris,epithet=microtis}
\newanimal{pd}{genus=Phyllostomus,epithet=discolor}
\newanimal{ef}{genus=Eptesicus,epithet=fuscus}
\newanimal{tb}{genus=Tadarida,epithet=brasiliensis}
\newanimal{gs}{genus=Glossophaga,epithet=soricina}

\newplant{me}{genus=Marcgravia,epithet=evenia}
\newplant{mh}{genus=Mucuna,epithet=holtonii}
\newplant{pp}{genus=Pachycereus,epithet=pringlei}

\newcommand{\flowchart}[1]{(Fig. \ref{fig:flowchart}#1)}

% *****************
%  Requirements:
% *****************
%
% - All pages sized consistently at 8.5 x 11 inches (US letter size).
% - PDF length <= 8 pages for full papers, <=2 pages for extended
%    abstracts (not including citations).
% - Abstract length <= 250 words.
% - No visible crop marks.
% - Images at no greater than 300 dpi, scaled at 100%.
% - Embedded open type fonts only.
% - All layers flattened.
% - No attachments.
% - All desired links active in the files.


% If your system does not generate letter format documents by default,
% you can use the following workflow:
% latex example
% bibtex example
% latex example ; latex example
% dvips -o example.ps -t letterSize example.dvi
% ps2pdf example.ps example.pdf


% For pdflatex users:
% The alifeconf style file loads the "graphicx" package, and
% this may lead some users of pdflatex to experience problems.
% These can be fixed by editing the alifeconf.sty file to specify:
% \usepackage[pdftex]{graphicx}
%   instead of
% \usepackage{graphicx}.
% The PDF output generated by pdflatex should match the required
% specifications and obviously the dvips and ps2pdf steps become
% unnecessary.


% Note:  Some laser printers have a serious problem printing TeX
% output. The use of ps type I fonts should avoid this problem.


\title{Sensorimotor behavior under informational constraints: a robotic model of prey localization in the bat \animal{mm}}
\author{Thinh Nguyen$^{1}$\and Dieter Vanderelst$^{2, *}$ \and Herbert Peremans$^{3}$\\
\mbox{}\\
$^1$ Electrical Engineering and Computer Science, University of Cincinnati\\
$^2$ Department of Biological Sciences, University of Cincinnati\\
$^3$Department of Engineering Management, University of Antwerp\\
* vanderdt@ucmail.uc.edu} % email of corresponding author

% For several authors from the same institution use the same number to
% refer to one address.
%
% If the names do not fit well on one line use
%         Author 1, Author 2 ... \\ {\Large\bf Author n} ...\\ ...
%
% If the title and author information do not fit in the area
% allocated, place \setlength\titlebox{<new height>} after the
% \documentclass line where <new height> is 2.25in

\pagenumbering{roman}

\begin{document}
\maketitle

\begin{abstract}
Action and perception are complicated by the fact that sensory information is often incomplete. Despite these informational constraints, animals interact intelligently and robustly with the world. Echolocating bats can be assumed to operate under substantial informational constraints as the information acquired by their sonar system is limited and sparse. Nevertheless, they fly swiftly through complex habitats and forage on the wing in complete darkness. Recently, the small tropical bat \animal[f]{mm} was documented to glean motionless prey from vegetation. Its ability to precisely localize insects has been interpreted as evidence for the bat acquiring a detailed acoustic image of the target. Given the limitations of the sonar system, it is unclear how this detailed image would be obtained. In this paper, we present a robotic model testing an alternative hypothesis. We propose that echo loudness (differences) contain sufficient information to localize a prey item perched on a leaf. Our robot can localize prey with a precision similar to that of the bat, relying only on echo loudness. We conclude that the documented behavior of \animal{mm} might not require a detailed acoustic image of the prey.
\end{abstract}

\section{Informational constraints}
%Echolocating bats have excellent spatial memory \citep{Barchi2013,Holland2007,Geva-Sagiv2015,Moehres1966}. They can navigate to salient locations like roosts, foraging grounds, and water sources \citep{Schnitzler2003,Helversen2003}. Displaced bats deprived of sight have been found to return to their roost successfully \citep{Stones1969,Williams1966}. Also, bats are highly maneuverable animals that can fly and forage along established routes in highly cluttered space \citep{Barchi2013,Petrites2009,VonHelversen2005}.

The extraordinary behavioral feats of echolocating bats have led researchers to speculate about the representations they infer from echoes. Some have suggested that the information acquired through sonar is highly detailed \citep{Simmons2012,Ulanovsky2008,Geipel2013}, containing information about object location and shape. Moreover, it has been suggested that their behaviors indicate that sonar enables bats to perceive complex scenes as integrated wholes \citep{Barchi2013}. And that the comprehensive acoustic images \citep{Clare2015} of the environment \citep{Lee2017} build from echo data compare well with those enabled by vision \citep{Moss2001}.

Despite the preceding interpretations of sonar-based behavior, bats can be assumed to operate under significant informational constraints, just like other animals. The world is full of intricate, ill-defined problems animals must solve to survive. Animals have to flexibly adjust their sensorimotor behavior to the moment-to-moment contingencies which arise in interaction with the environment \citep[e.g.,][]{Brooks1999,Beer1990,Pfeifer2006,DiPaolo2017}. Critically, action and perception have to proceed in the face of substantial informational deficiency \citep{Brooks1990,Sterling2015}. Indeed, animals can not be assumed to have access to complete knowledge about (the dynamics of) their bodies or the environment. This informational constraint encumbers action-perception in all animals \citep{Marks1988,Sterling2015}, including mammals, and is one of the primary reasons direct interaction with the world is challenging \citep{Alterovitz2016}. 

The view that bat echolocation yields highly detailed information about the environment probably has its origin in experimental findings showing that bats can locate single targets (typically, a small sphere) with high accuracy \citep[e.g.,][]{Simmons1983,Lawrence1982,Wotton2000}. However, the ability to locate and identify isolated objects under favorable conditions does not necessarily imply that bats can reconstruct the layout of multiple complex reflectors (e.g., trees and shrubs) under natural conditions. In contrast to the artificial targets used in experiments, most natural objects, including vegetation and human-made objects, typically return many overlapping echoes \citep{Yovel2011,Yovel2009,Vanderelst2016,Warnecke2018}.

Over the past few years, we have shown that for a range of sonar-based activities, bats do not require an interpretation of echoes in terms of identifiable objects in the environment or acoustic images  \citep[e.g.,][]{Vanderelst2015,Vanderelst2016,Steckel2013,Mansour2019a,Achutha2021}. We maintain that bat sonar is inherently limited in the amount of information it can extract about natural environments' structure. The presumed ability of bats to reconstruct acoustic images from echoes is questionable due to their small field of view \citep{Surlykke2009,Jakobsen2012,Jakobsen2013} and low update rate \citep{Holderied2006,Seibert2013}. However, the most severe barrier stems from the temporal resolution of the sonar system \citep{Simmons1989,Wiegrebe1996,Surlykke1996}.  It follows that bats probably cannot extract much information about the 3D layout of the environment most of the time. 

\section{Gleaning prey from substrate}

Recently, \citet{Geipel2013, Geipel2019} reported on prey capture behavior in the small neotropical bat \animal[f]{mm} (see fig. \ref{fig:setup}a). This bat hunts for insects perched on the leaves of understory vegetation in the rainforest. This behavior is remarkable as echoes from the leaves overlap with any echoes returning from potential prey \citep{Jones2013}. The overlapping echoes should make it difficult to detect and localize the prey. 

\citet{Geipel2013, Geipel2019} presented some hypotheses about how \animal{mm} can hunt for motionless prey resting on leaves. Ensonification experiments confirmed that leaves with an insect on them return louder echoes than unoccupied leaves, especially when approached under an oblique angle \citep{Geipel2019}. While this might explain how \animal{mm} detects prey (or, rather, leaves occupied by prey), it does not explain how the bat localizes and homes in on the target.

\citet{Geipel2013} suggested that the bat's ability to localize the insect indicates that they form acoustic images from temporal and spectral
echo information \citep[see also][]{Geipel2007,Jones2013}. Likewise, \citet{Simon2014} interpreted the findings of \citet{Geipel2013} as evidence for fine echo-acoustic resolution. If the bat can form a 2D or 3D representation of the leave, the insect, or combination thereof, it could use this to approach the insect. However, we think it is unlikely that \animal{mm} can infer the insect's 3D geometry and leaf from echoes for the reasons outlined above. In this paper, we test the alternative hypothesis that a strategy based on echo amplitude allows \animal{mm} to approach prey perched on a leaf. We test this hypothesis using a robotic model operating under conditions mimicking the experiments presented by \citet{Geipel2013, Geipel2019}. 

\section{Methods}

\subsection{Robot and setup}

We equipped an ST Robotics R12 robot arm with a sonar head (figure \ref{fig:setup}). The robot arm was mounted on a linear track with a length of 3 meters. The sonar head consisted of two Maxbotix MB1360 sonar sensors. These sensors emit echo pulses with a frequency of 42 kHz. The sensors return the envelope of the resulting echo waveform as an analog signal. Pulsing the sensors and reading out their responses was done through an onboard microcontroller. The microcontroller communicated with a host computer over Bluetooth. The sonar head was powered using a 9V battery.

\begin{figure*}[htb]
	\centering
	\includegraphics[width=0.75\linewidth]{plots/setup}
	\caption{Setup for the robotic experiments. \textbf{(a)} Example of experiments conducted by \citet{Geipel2019} documenting the prey gleaning behavior of \animal{mm}. In these experiments, the bat is presented with a dragonfly on a leaf. The bat's prey gleaning behavior is documented using high-speed infrared cameras. \textbf{(b)} Detail of the artificial dragonfly and leaf used in the robotic experiments. \textbf{(c)} General view of the setup showing the robot arm and the artificial leaf, with the dragonfly. \textbf{(d)} Detail of the sonar head. \textbf{(e)} Coordinate frames used in the paper. Axes $X_s$, $Y_s$ and $Z_s$ form a coordinate frame attached to the sonar head. The axis $X_s$ is perpendicular to the sensor surface. $X_d$ (normal to the leaf), $Y_d$ (parallel to leaf), and $Z_d$ denote a coordinate frame attached to the dragonfly. The $X_d$ axis is perpendicular to the leaf surface. The robot's strategy was designed to try and regulate the azimuth $\azimuth$ and elevation $\elevation$ angle of the dragonfly, in the $X_s$, $Y_s$ (pointing away from the viewer into the page) and $Z_s$ coordinate frame, to zero. If both azimuth and elevation are zero, the sensor is focusing on the dragonfly. $X_w$, $Y_w$ (not drawn), and $Z_w$ are the world coordinate frame.}
	\label{fig:setup}
\end{figure*}

The sonar sensors were mounted with their acoustic axes offset by 30 degrees in azimuth. This modelled the orientation of the acoustic axes of bats ears, as observed in their \hrtfs\ \citep[e.g.,][]{DeMey2008}, including \animal{mm} \citep{Vanderelst2010,Firzlaff2004}. To collect an echo measurement, each of the two sonar sensors was triggered in succession, with a short time interval to avoid cross-talk between the sensors.

In experiments observing \animal{mm} capturing prey perched on leaves, the bat used short ($\sim 0.2$ ms) ultrasonic calls with a peak frequency around 100 kHz \citep{Geipel2013}.  Our sonar sensors operate at a frequency of 42 kHz. Therefore, as detailed below, we scaled all elements in our setup by a factor 2 to obtain consistency between the size of the artificial leaf and insect and the wavelengths used. This is a conservative scaling factor as $100\ \mbox{kHz} / 42\ \mbox{kHz} > 2$. Moreover, \animal{mm} uses frequencies well above 100 kHz, and these might be most informative in localizing prey \citep{Geipel2019}.

\citet{Geipel2019} report on experiments documenting the prey capture behavior of \animal{mm}. They presented the bat with prey items perched on an artificial heart-shaped leaf measuring 21 cm in height and 11.6 cm in width (Fig. \ref{fig:setup}a). We constructed a heart-shaped leaf measuring 42 ($2 \times 21$ cm) cm by 22 cm ($2 \times 11$ cm) out of heavyweight paper. This artificial leaf's shape was obtained by scaling a picture of a philodendron leaf (figure \ref{fig:setup}b).

The dragonfly used by \citet{Geipel2019} had a body length of about 4.5 cm. We created a 3D model of a dragonfly body by tracing a drawing of the dragonfly \animal{ca} \citep[][plate V]{Dragonflies1890}. The model was scaled to have a body length of 9 cm ($2 \times 4.5$ cm). This resulted in a wingspan of about 12 cm. The body was 3D printed using a Form2 resin printer. The wings were traced on paper, cut out, and glued onto the 3D-printed body. The dragonfly used by \citet{Geipel2019} protruded 2 cm from the surface of the leaf. Therefore, we provided the 3D model of the dragonfly with a foot that raised it 4 cm from the cardboard leaf's surface. The foot of the artificial dragonfly was equipped with a small metal disk. This enabled us to position the model on the leaf employing a magnet. The setup, and the coordinate frames, are illustrated in figure \ref{fig:setup}.

Our artificial leaf and dragonfly were constructed of different materials than their natural counterparts. However, the intensity reflection coefficient depends on the difference in acoustic impedance between air and any reflectors \citep{Alkins2014}. Air has a very low acoustic impedance compared to solid materials. This results in echoes depending mainly on the shape of the reflector and only minimally on its material. Therefore, we expect echoes in our setup to depend on the leaf and artificial dragonfly shape rather than any differences in material.

\subsection{Dragonfly approach strategy}

To test whether the bat \animal{mm} could approach a prey item perched on a leaf without a detailed acoustic image of the prey \citep{Geipel2013}, we implemented a simple approach strategy in the robot. The algorithm is illustrated in figure \ref{fig:flowchart}. In the text below, we refer to the labels in figure \ref{fig:flowchart}.

\begin{figure}[tb]
	\centering
	\includegraphics[width=1\linewidth]{plots/flow_diagram}
	\caption{Flowchart summarizing the procedure we implemented for the robot to localize the dragonfly on the leaf. The labels are referred in the main text.}
	\label{fig:flowchart}
\end{figure}

Each trial started with randomly positioning the sonar head in a volume bounded by the following coordinates. The subscript $w$ refers to world coordinates, while the subscript $d$ refers to a coordinate frame attached to the dragonfly, see fig. \ref{fig:setup}e. The dragonflies world coordinates were $[x_w = 1000, y_w = 0, z_w = 500]$. All coordinates are in millimeters.

\begin{align} 
x_w = 300\ &\mbox{or}\ x_d = 700\\ 
y_w \in [-300, -100]\ &\mbox{or}\ y_d \in [300, 100] \\
z_w \in [100, 300]\ &\mbox{or}\ z_d \in [-400, -200]
\end{align}

This positioned the sonar head 400 to 200 mm below the dragonfly and up to 300 mm left of the midline. The distance along the $X_d$ axis was 700 mm. Also, the sonar sensor assumed a yaw $\yaw_s$ (rotation around $Z_s$, positive $\yaw_s$ toward the $Y_s$ axis) and pitch $\pitch_s$ (rotation around $Y_s$, positive $\pitch_s$ away from the $Z_s$ axis) sampled from the following intervals,

\begin{equation}
\yaw_s \in [20^{\circ}, 45^{\circ}]; \pitch_s \in [-10^{\circ}, -15^{\circ}]
\end{equation}

This rotated the sensor such that it looked somewhat upwards and towards the midline ($X_w$ axis, fig. \ref{fig:setup}e). This ensured that the robot sensor's initial position was suitable to pick up echoes from the dragonfly and leaf in at least one ear (sonar sensor).

By positioning the robot as just described at the start of each trial, we model the behavior of \animal{mm}. Indeed, the animal was found to typically approach leaves occupied with insects from below \citep{Geberl2013,Geipel2019}. Therefore, our robotic model can be interpreted as mimicking a bat that has decided that a leaf is worth approaching but has not yet precisely localized it (or, more importantly, the insect upon it). Modeling localizing the prey as a separate stage in the hunting behavior of \animal{mm} corresponds to this animal's observed behavior. \citet{Geipel2007} described how an initial approach of an occupied leaf is typically followed by a scanning behavior that leads to localization of the prey.

After moving the sonar head to its initial location, the approach behavior began. The algorithm we implemented is depicted in figure \ref{fig:flowchart}. In brief, the algorithm tried to localize the dragonfly by (1) minimizing the echo amplitude difference between the two ears and (2) maximizing the echo amplitude summed across the two ears. Minimizing the difference across ears was used to nullify the dragonfly's azimuth $\azimuth_d$ in the sensor's coordinate frame. Maximizing the amplitude across ears aimed at nullifying the dragonfly's elevation $\elevation_d$. If both  $\azimuth_d$ and $\elevation_d$ are zero, the sonar sensor is pointing directly at the dragonfly. In the next paragraphs, we describe the algorithm in more detail and specify how it corresponds to the information available to bats, and \animal{mm} in particular.

The algorithm started by taking a single measurement, pulsing both sensors in succession \flowchart{a}. The combined echo of the leaf and dragonfly was assumed to start at the first above threshold, after the pick up of the emission was removed. Extracting samples in a $\sim$0.5 ms window in the left $s_{l,i}$ and the right ear $s_{r,i}$ envelope signals, the interaural level difference IID was calculated as follows,

\begin{equation}
    \mbox{IID} = 10 \times log10 \frac{\sum_i s_{l,i}}{\sum_i s_{r,i}}
\end{equation}

The IID value was used to estimate the azimuth position $\hat{\azimuth}_d$ of the dragonfly \flowchart{b}. This was done using an empirically determined function that related IID to the azimuth position of a reference object, i.e., a 10 cm diameter cardboard cylinder. Using these measurements allowed us to compensate for differences in sensitivity between the left and the right sonar sensor in the sonar head. If the absolute value of the estimated azimuth position $\hat{\azimuth}_d$ of the dragonfly was smaller than 5 degrees \flowchart{c}, the algorithm attempted to nullify elevation $\elevation_d$ \flowchart{d-e}. 

Bats' \hrtfs\ filter incoming echoes, imputing information about the elevation direction of the echo \citep[e.g.,][and references therein]{Reijniers2007,Reijniers2010,Vanderelst2010}. Our robot can not use this frequency-dependent filtering as it only uses a very narrow range of frequencies around 42 kHz. However, one source of elevation information in the HRTFs of bats is yielded by the frequency-dependent change of the maximum sensitivity axis. In other words, the main acoustic axis of bats' sonar systems sweeps across azimuth and elevation as a function of frequency. This is also true for the bat \animal{mm}. In previous work, we simulated the \hrtf\ of this animal \citep{Vanderelst2010}.

In the frequency range of the calls used while hunting for prey on leaves, the main axis sweeps across about 30 degrees in elevation. Therefore, to localize the dragonfly in elevation, we performed five measurements, equally spaced from -15 to 15 degrees around the sensor's current pitch angle $\pitch_s$ \flowchart{d}. For each pitch position, the samples of the echo signals in the left and the right ear were summed, giving the echo amplitude $L$ across the two ears,

\begin{equation}
    L_{\pitch_s} = \sum_i s_{l,i,\pi_s} + \sum_i s_{r,i,\pitch_s}
\end{equation}

The estimated elevation of the dragonfly $\hat{\elevation}_d$ was determined as equal to the pitch angle $\pitch_s$ yielding the largest value $L_{\pitch}$ \flowchart{e}.

After estimating the azimuth and elevation positions of the dragonfly $\hat{\azimuth}_d$ and $\hat{\elevation}_d$, the following corrections to the yaw and pitch of the sonar were applied \flowchart{f},

\begin{align}
    \Delta \yaw_s &= \hat{\alpha}_d \times 0.75\\
    \Delta \pitch_s &= \hat{\epsilon}_d \times 0.75
\end{align}

which yield a simple proportional controller. After applying these corrections to the sonar head's orientation, the sensor was translated by 20 mm, in the direction the sonar sensor was pointing in \flowchart{g}, i.e., along the $X_s$ axis. This displacement was based on the finding that \animal{mm} vocalizes about 50 times per second \citep{Geipel2013}. Moreover, we reanalyzed the data collected by \citet{Geipel2019}.  We found that during its approach of a leaf, the bat maintains a speed of about 0.5 m/s. Therefore, the bat moves about 1 cm between calls. In our scaled setup, this translates to 20 mm.

The procedure outline above was repeated until the sensor was 400 mm from the dragonfly \flowchart{h}. Considering the scaling of the setup, this would correspond to a distance of 20 cm for the bat. We stopped the approach at this point because we found that, when closer than about 40 cm from the leaf, the echoes tended to saturate the receiver. Thereby no longer enabling us to use echo amplitude (differences) to localize the dragonfly. 

The bat should be able to approach the leaf to about 7 cm without overlap between pulse and echo as it uses extremely short emissions \citep[0.2 ms,][]{Geipel2013}. Moreover, bats regulate the emission amplitude to ensure that returning echoes are neither too loud or too faint. Our sonar system could not mimic these very brief pulses, even after taking into account the scaling of the setup. Moreover, the sensors do not allow to reduce the emission amplitude to avoid saturation of the receivers.

\subsection{Conditions and trials}

We tested the algorithm outlined above in two conditions. In one condition, we placed the dragonfly in the middle of the leaf, at the same relative position as done by \citet{Geipel2019}. In a second condition, we moved the dragonfly 90 mm towards the edge of the leaf. By testing the robot's approach performance towards a dragonfly placed at different positions, we could assess whether the robot is localizing the dragonfly instead of approaching the center of the leaf. For each condition, we ran 25 trials.

\section{Results}

The results\footnote{All data included in this paper are available online at \url{https://doi.org/10.6084/m9.figshare.14582493.v1}} for the dragonfly in the middle of the leaf are presented graphically in figure \ref{fig:results_middle}. Numerical results are listed in table \ref{tab:results_middle}. 

\begin{figure*}[tb]
	\centering
	\includegraphics[width=1\linewidth]{plots/results_middle.pdf}
	\caption{Results for the dragonfly positioned in the middle of the leaf. \textbf{(a)} Evolution of azimuth error $\azimuth_d$ as function of trial step. \textbf{(a)} Similar as panel (a), but for elevation error $\elevation_d$. To compensate for the scaling of the setup, the data were divided by a factor 2 before plotting. \textbf{(c)} and \textbf{(d)} Trajectories of the sensor in world coordinates, $X_w$, $Y_w$, and $Z_w$. In these panels, the black dot indicates the position of the dragonfly. The differently colored lines and markers correspond to different trials. \textbf{(e)} The Cartesian positions showing the sensor's focus of attention projected onto the plane of the dragonfly at the end of each trial. The colors of trials correspond across panels.}
	\label{fig:results_middle}
\end{figure*}

\begin{table}[tb]
\centering
\input{plots/results_middle}
\caption{Numerical results for the condition with the dragonfly in the middle of the leaf. The table lists the azimuth $\azimuth_d$ and elevation $\elevation_d$ of the dragonfly at the end of the approach sequences (in degrees, corrected for the applied scaling). The table also lists statistics on the final error in Cartesian coordinates. These are the $y_d$ and $z_d$ coordinates the sensor pointed to at the end of the approach. This indicates the point at which the $X$-axis of the sensor,$X_s$, would intersect the leaf (in millimeters, corrected for the applied scaling).}
\label{tab:results_middle}
\end{table}

Figure \ref{fig:results_middle} depicts the evolution of the dragonfly's position as a function of time (iteration) for each trial (fig. \ref{fig:results_middle}a, b). In particular, we plot the azimuth $\azimuth_d$ and elevation angle $\elevation_d$. Recall that, if $\azimuth_d = 0$ and $\elevation_d = 0$, the sensor fixated the dragonfly, particularly the center of the thorax (see figure \ref{fig:setup} for the coordinate frames employed in this paper). Figure \ref{fig:results_middle}c, d shows the trajectories of the robot in world coordinates $x_w$, $y_w$ and $z_w$. Finally, figure \ref{fig:results_middle}e shows the sensor's final orientation, projected in the dragonfly's plane. This is, it shows the position where the sensor would strike the dragonfly or leaf if it were to move forward in the direction of its final yaw $\yaw_s$ and pitch $\pitch_s$ orientation.

Figure \ref{fig:results_middle} shows that the robot could locate the dragonfly on most runs. The azimuth error $\azimuth_d$ and elevation error $\elevation_d$ tended to decrease over time. At the end of most runs, the sensor pointed towards the dragonfly's body, i.e. the azimuth errors $\azimuth_d$ and elevation errors $\elevation_d$ were small (Median values: \MiddleAz$^{\circ}$, \MiddleEl$^{\circ}$, figure \ref{fig:results_middle} and table \ref{tab:results_middle}). In Cartesian coordinates, the median horizontal and vertical errors were \MiddleDy\ mm and \MiddleDz\ mm, respectively. However, on several runs, the robot did miss the dragonfly. In these cases, the robot seems mostly to fixate the edge of the leaf. The localization in elevation was less precise than localization in azimuth ($\rho_{\elevation_d} > \rho_{\azimuth_d}$, Levene's test for equality of variance, \MiddleTest).

The results for the dragonfly in the middle of the leaf do not necessarily demonstrate that the robot located the prey. Alternatively, the robot could have been focusing on the center of the leaf. Therefore, we also ran trials in which the dragonfly was shifted towards the edge of the leaf. The results for these runs are plotted in figure \ref{fig:results_side}. Numerical results are listed in table \ref{tab:results_side}.

As for the results with the dragonfly in the center of the leaf, the robot could locate the dragonfly in most runs. Most importantly, the final orientations of the sensor are appropriately shifted with respect to the leaf. This shows that the robot did not just locate the center of the leaf. Instead, it was able to find the dragonfly on the leaf. As in the previous condition, the azimuth errors $\azimuth_d$ and elevation errors $\elevation_d$ at the end of the trials tended to be small (Median values: \SideAz$^{\circ}$, \SideEl$^{\circ}$, figure \ref{fig:results_side} and table \ref{tab:results_middle}). In Cartesian coordinates, the median horizontal and vertical errors were \SideDy\ mm and \SideDz\ mm, respectively. Again, the localization in elevation was less precise than localization in azimuth ($\rho_{\elevation_d} > \rho_{\azimuth_d}$, Levene's test for equality of variance, \SideTest). Moreover, we also observed a tendency to focus the edge of the leaf, as in the previous condition.

\begin{figure*}[tb]
	\centering
	\includegraphics[width=1\linewidth]{plots/results_side.pdf}
	\caption{Similar as figure \ref{fig:results_middle}, but for the trials with the the dragonfly positioned on the edge of the leaf.}
	\label{fig:results_side}
\end{figure*}

\begin{table}[tb]
\centering
\input{plots/results_side}
\caption{Numerical results for trials in which the dragonfly was placed at the edge of the leaf. Data similar to table \ref{tab:results_middle}.}.
\label{tab:results_side}
\end{table}

\subsection{Bat data}

To the best of our knowledge, the prey localization precision of \animal{mm} has not been evaluated formally. \citet{Geipel2007} and \citet{Geipel2013} mentioned that \animal{mm} almost always bit the dragonfly in the thorax. However, the video data collected in a subsequent experiment \citep{Geipel2019} show that the bat often initially misses the dragonfly, misdirecting its bite. This is then followed by a re-adjustment of the bite to the thorax, body, or wings. Our setup more closely mimics the experiments reported by \cite{Geipel2019} (compare fig. \ref{fig:setup}a and b). Most importantly, the leaf and the prey's relative size were markedly different in the experiments of \citet{Geipel2007} and \citet{Geipel2013}. Therefore, we re-assessed the high-speed video collected in this more recent study to obtain data, allowing us to compare the bat and the robot's behavior.

Figure \ref{fig:bat_captures} shows stills taken from different trials in the study by \citet{Geipel2019}. In each panel, we show the bat just before making contact with the leaf or dragonfly. In each of these images, the bat strikes the leaf, initially misdirecting its bite below, above or to the side of the dragonfly. In the videos, the bat subsequently adjusts its bite and manages to capture the dragonfly.

\begin{figure*}[htb]
	\centering
	\includegraphics[width=1\linewidth]{plots/bat_captures.pdf}
	\caption{Stills from high-speed camera footage documenting the prey gleaning behavior of \animal{mm} \citep{Geipel2019}. Because the dragonfly is hard to see in the images, the white arrows indicate the extent of its body. In all trials, the bat initially made contact with the leaf instead of the dragonfly.}
	\label{fig:bat_captures}
\end{figure*}

\section{Discussion}

From the data presented in figures \ref{fig:results_middle}, \ref{fig:results_side} and tables \ref{tab:results_middle}, \ref{tab:results_side}, we conclude that a simple intensity (difference) driven approach strategy is sufficient to locate  a dragonfly on a leaf. Trials with the dragonfly at the edge of the leaf confirmed that the robot located the prey instead of the leaf. The errors are comparable to those for the first condition.

From the video stills presented in figure \ref{fig:bat_captures}, we infer that the robot's errors are in the same range as the bat's precision of attack. Indeed, in none of the trials does the bat bite the dragonfly in the thorax. Instead, its initial bite misses the dragonfly in all three examples.

In both conditions, results show a cluster of 'attacks' directed at the artificial leaf's edge. To the best of our knowledge, these systematic errors have not been observed in the bat. However, it should be pointed out that the available data for the bat is minimal. Perhaps, additional data on the bat might reveal recurring types of errors. These might also further inform our modeling of \animal{mm}.

In interpreting the robotic results, several disadvantages of our robot with respect to the bat should be kept in mind. The most important is the fact that \animal{mm} uses broadband calls spanning at least 50-140 kHz (Wavelengths: 7-2.5 mm) \cite{Geipel2019}. Even using only echo intensity to guide the approach, this broad range of frequencies would be advantageous. The amplitude of broadband echoes is much less sensitive to destructive and constructive interference because of the range of constituent wavelengths. Destructive interference at one wavelength is likely to be compensated by constructive interference at another. In contrast, our robot used a narrowband call centered around 40 kHz, resulting in augmented noise on the echo intensities' estimates. 

A second disadvantage was that our robot did not have access to a real \hrtf. In the absence of an \hrtf\ providing elevation localization cues, we approximated this by gathering echoes for five different pitch steps. Perhaps, this approximation resulted in increased elevation errors. In both conditions, the elevation errors $\elevation_d$ tended to vary more than the azimuth errors $\azimuth_d$, which could suggest that our approximation might be improved upon.

Finding that a simple intensity-based algorithm can mimic the prey localization observed in  \animal{mm} suggests that no detailed acoustic image \citep{Geipel2013} is required to find the prey on a leaf. In subsequent research, we aim to increase our robot's sensor abilities to match the bat's sonar sensor better. In addition, we will test the robot in more challenging and diverse conditions. For example, it would be interesting to assess whether intensity-driven approaches allow selecting leaves occupied by insects among many leaves.

\footnotesize
\bibliographystyle{apalike}
\bibliography{references} % replace by the name of your .bib file

\end{document}
